 \documentclass[11pt]{article}
\usepackage{amsmath,amsfonts}
\usepackage{graphicx}
%\usepackage{xcolor}
%\definecolor{gray}{rgb}{.75,.75,.75}
%\usepackage[landscape]{geometry}
%\usepackage{hyperref}
%\usepackage[bookmarks=true, pdfborder={0 0 .5}, urlbordercolor=gray]{hyperref}

\renewcommand{\thefootnote}{\fnsymbol{footnote}}


\addtolength{\topmargin}{-.75in}
\addtolength{\textheight}{1.5in}
\addtolength{\oddsidemargin}{-.35in}
\addtolength{\textwidth}{.7in}
\pagestyle{empty}

\newcommand{\ds}{\displaystyle}
\newcommand{\ts}{\textstyle}

\newcommand{\Z}{\mathbb{Z}}
\newcommand{\R}{\mathbb{R}}
\newcommand{\C}{\mathbb{C}}
\newcommand{\EE}{\mathcal{E}}
\newcommand{\tZ}{\tilde{\Z}}

\newcommand{\Sym}{\mathrm{Sym}}

\renewcommand{\circle}[1]{{\bigcirc\hspace{-.75em}#1}}

\begin{document}

\footnote[0]{Notes by Peter Magyar \ {\tt magyar@math.msu.edu}}
\vspace{-1em}

\mbox{}\vspace{1em}

\centerline{\bf Math 133\hspace{2.5em} \hfill{\Large\bf
Polar Coordinates}\hfill 
Stewart \S10.3/I,II}
\vspace{1em}

\noindent 
{\bf Points in polar coordinates.} The first and greatest achievement of modern mathematics
was Descartes' description of geometric objects by numbers, using a {\it system of coordinates}.
In the simplest example, {\it Cartesian} or {\it rectangular} coordinates on the plane 
locate a point $P$ in terms of two coordinate measurements $x$ and $y$:
how far over and how far up the point is, moving parallel to the marked axes.
We loosely say that $P$ ``is'' the pair $(x,y)$, because the coordinates 
tell how to get there from the origin.
The name $P$ is like identifying a house as ``the Jones place'', 
whereas the coordinates are like saying ``the third house to the 
right on the second street down''.

In this section, we learn how to locate the point $P$ using a different pair of measurements,
the {\it polar coordinates} $(r,\theta)$. 
The {\it radius} $r$ is the distance from the origin.
The {\it angle} $\theta$ is measured couterclockwise in radians, 
from the positive $x$-axis ray to the ray from orgin through the point $P$.
This is like pointing to ``the house 500 yards in {\it that} direction''.
%, a description more useful in the countryside than in the city.
\\[.7em]
\centerline{
\includegraphics[width=1.6in]
{10-3-Polar-Point.png} 
}
\vspace{-1em}

Unlike rectangular coordinates, the polar coordinates of a point 
are multivalent, having many equivalent versions because of the ambiguity of angles.
For example, the point $(x,y)=(0,1)$ on the positive $y$-axis 
corresponds to $(r,\theta)=(1,\frac\pi2)$, 
where $\theta=\frac\pi2$ means a $\frac14$ turn counterclockwise 
from the positive $x$-axis.
However, we could equally well get to this point by a $\frac34$ turn clockwise,
giving $(r,\theta)=(1,-\frac{3\pi}2)$. 
In fact, we could get to the point by $1\frac14$ turns counterclockwise,
$1\frac34$ clockwise, etc.
In general, we must consider all angles that differ by a multiple of a full turn $2\pi$ as the same,
meaning they define the same point:
$$
(r,\theta) = (r,\,\theta{+}2n\pi) \quad\text{for any integer $n$.}
$$
It is also useful to allow negative radius: $(-r,\theta)$ means to move
out along the line at angle $\theta$, but in the opposite direction from the positive ray,
along the ray $\theta\pm\pi$; thus:
$$
(-r, \theta) = (r,\,\theta\pm\pi)\,.
$$
There is even more ambiguity for the origin $(x,y)=(0,0)$,  
which can be written as $(r,\theta)=(0,\theta)$ for any angle at all.

Both types of coordinates completely locate a point, 
so given either $(x,y)$ or $(r,\theta)$,
we can find the other by simple trigonometric formulas:
$$\begin{array}{lcl}
\text{Given $(r,\theta)$}
&\quad\Longrightarrow\quad&
\text{find $(x,y)$ with}\ 
\left\{\begin{array}{rcl}
x &=& r\cos(\theta)
\\[.3em]
y &=& r\sin(\theta)\,.
\end{array}\right.
\\[2em]
\text{Given $(x,y)$}
&\quad\Longrightarrow\quad&
\text{find $(r,\theta)$ with}\ 
\left\{\begin{array}{rcl}
r &=& \sqrt{x^2+y^2}
\\[.3em]
\theta &=& \arctan(\tfrac yx)\,.
\end{array}\right.
\end{array}$$
Here we get $\theta$ from the defining formula $\tan(\theta)=\frac yx$,
and we could equally well use $\sin(\theta) =\frac{y}{\sqrt{x^2{+}y^2}}$, etc.,
always remembering we can change $\theta$ to
$\theta{+}2n\pi$.
Also, since $-\frac\pi2<\arctan(\theta)<\frac\pi2$,
we must define $\arctan(\infty)=\frac\pi2$, 
$\arctan(-\infty)=-\frac\pi2$;
and we must adjust the angle by $\pm\pi$ if the point lies left of the $y$-axis.\footnote{
Summarizing:
$$
\theta 
\ =\ 
\arcsin\!\left(\!\frac{y}{\sqrt{x^2{+}y^2}}\!\right)
\ =\ 
\arccos\!\left(\!\frac{x}{\sqrt{x^2{+}y^2}}\!\right)
\ =\ 
\left\{\!\!\begin{array}{cl}
\arctan\left(\frac yx\right) &\text{if $x\geq 0$}\\[.5em]
\arctan\left(\frac yx\right) +\mathrm{sgn}(y)\pi &\text{if $x< 0$}
\end{array}\right.
$$
The function after the last equality is called {\tt atan2(y,x)} in computer languages.
}
\\[1em]
{\bf Curves in polar coordinates.}
Any geometric object in the plane is a set (collection) of points, 
so we can describe it by a set of coordinate pairs.
For example,  the unit circle $C$ is the set of all points at distance 1 from the origin;\footnote{ 
There is no separate curve ``connecting'' the points: the curve {\it is} just all the points.}
the coordinates of these points form the set of all pairs $(x,y)$ which satisfy the Pythagorean equation $x^2 + y^2=1$:
$$
C \ = \ \{(x,y)\text{ such that } x^2+y^2=1\}\,.
$$
Again, the equality of these sets is meant loosely: a pair of numbers
like $(x,y)=(\frac35,\frac45)$ is not literally a geometric point  on the circle, 
but it {\it identifies} a point by means of the rectangular coordinate system.
Now, polar coordinates are specially adapted to describe round,
turny shapes centered at the origin,
and they make the equation of the circle as simple as possible:
$$
C \ = \ \{(r,\theta)\text{ such that } r=1\}\,.
$$
We also have the parametric forms $(r(t),\theta(t))=(1,t)$ and
$(x(t),y(t))= (1\cos(t),1\sin(t))$.
\\[.5em]
{\sc example:} The line $x+y=1$ is not at all circular or centered at the origin, and its equation becomes complicated 
in polar coordinates:
$$
x+y=1
\ \ \Longleftrightarrow\ \
r\cos(\theta) + r\sin(\theta) \ =\  1
$$
$$
\ \ \Longleftrightarrow\ \
r \ =\ \frac{1}{\cos(\theta){+}\sin(\theta)}
\ =\ 
\tfrac1{\sqrt2}\sec(\theta{-}\tfrac\pi4)\,.
$$
The last equality follows from the identity
$\cos(\theta{-}\frac{\pi}4) = \cos(\theta)\cos(-\frac\pi4) - \sin(\theta)\sin(-\frac\pi4)
=\frac1{\sqrt2}(\cos(\theta)+\sin(\theta))$.
Parametrically, $(r(t),\theta(t))=(\frac{1}{\sqrt 2}\sec(t{-}\frac{\pi}4),t)$, $t\in (-\frac{\pi}4,\frac{3\pi}4)$;
also $(x(t),y(t))=\frac{1}{\sqrt 2}\sec(t{-}\frac{\pi}4)(\cos(t), \sin(t))$,
a trajectory with constant {\it angular} velocity.

Similar reasoning gives the polar form of a general linear equation.
For $ax + by=0$, we get $\theta=\alpha+\frac\pi2$ for the constant angle $\alpha=\arctan(\frac ba)$.
For constant term $c\neq 0$:
$$
ax + by = c
\quad\Longleftrightarrow\quad
r \ = \ \frac{c}{a\cos(\theta)+b\sin(\theta)}
\ =\
\frac{c}{\sqrt{a^2+b^2}}\, \sec(\theta{-}\alpha)\,.
$$

\newpage

\noindent
{\sc example:} Consider Archimedes' spiral, 
the shape of the groove on an old vinyl record (solid blue line).
\\[.0em]
\centerline{
\includegraphics[width=2.5in]
{10-3-Spiral.png} 
}
\vspace{-.5em}

\noindent
This is defined by a point moving steadily outward
as it turns around the origin: in parametric polar coordinates,
$(r(t),\theta(t))= (t,t)$ for $t\geq 0$,
meaning at time $t$ the radius and angle are both $t$.
Converting into rectangular coordinates:
$$
(x(t),y(t))\ =\
(r(t)\cos \theta(t), r(t)\sin\theta(t))
\ =\
(t\cos(t), t\sin(t))\,, \quad t\geq 0.
$$
Deparametrizing gives the $r\theta$ and $xy$-equations:
$$
r \ =\ \theta + 2n\pi \ \ \text{for integer $n$}
\quad\Longrightarrow\quad
\sqrt{x^2+y^2} \ =\ \arctan\!\left(\frac yx\right)+2\pi n
$$
\\[-2.5em]
$$
\quad\Longrightarrow\quad
y \ = \ x\,\tan\!\sqrt{x^2+y^2}\,.
$$
\\[-1em]
For example, we can tell the points $(x,y) = (2n\pi, 0)$ are on the spiral, because
$0 = 2n\pi\tan\sqrt{(2n\pi)^2+0^2}$.
Actually, this $xy$-equation defines the spiral together with
its $180^\circ$ rotation (dashed red line);
note that this is different from the continuation of the $r\theta$-equation
$r=\theta$ for $\theta<0$, which is the left-to-right reflection of the spiral.
\\[1em]
{\bf Sketching polar graphs.}  Remember that a function $f$ 
is just a rule taking input numbers to output numbers.
It does not care what letters we use for inputs and outputs,
or how we interpret those letters geometrically.
We usually illustrate the function by drawing its rectangular graph
$y=f(x)$, in which $f$ controls the height $y$ above each point on the
$x$-axis. But another way to illustrate this function is the polar graph
$r=f(\theta)$, in which $f$ controls the radius $r$ along each 
ray $\theta$.

You can sketch the polar graph $r=f(\theta)$ by plotting points, 
just as for a rectangular graph. For example, consider the polar curve:
$$
r = \sin(\theta)\,.
$$
Imagine the plane as a field, with you standing at the origin. 
Look along the positive $x$-axis in direction $\theta=0$,
and draw a point at radius $r=\sin(0)=0$, namely the origin itself.  
As you increase $\theta>0$, turning slowly to the left, 
you increase the radius as $\sin(\theta)$ increases.  
The radius tops out at 1 when $\theta=\frac\pi2$ along the positive $y$-axis;
and as you continue to turn, the point comes back in to the origin when $\theta=\pi$.
After that, as you turn toward negative $y$ directions
the radius becomes negative, so you draw points behind you, 
in fact retracing the original curve.
\\[-1em]

\centerline{
\includegraphics[width=2.5in]
{10-3-Sin-Theta.png} 
}
\vspace{0em}

\noindent
This is a computer plot to turn the qualitative story above into an accurate graph. 
But you really could do this by hand, by plotting $r$ for some standard $\theta$:
$$
\begin{array}{|c||c|c|c|c|c|c|c|c|c|} 
\hline &&&&&&&&& \\[-.8em]
\text{deg}&\ 0^{\circ} & 30^{\circ}\!\! & 45^{\circ}\!\! & 60^{\circ}\!\! & 90^{\circ}\!\! & \!\!120^{\circ}\!\!\! &\!\! 135^{\circ}\!\!\! &\!\! 150^{\circ}\!\!\! &\!\!  180^{\circ}\!\!\! 
 \\[.3em] \hline&&&&&&&&& \\[-.8em]
\theta & 0 & \frac{\pi}6 & \frac{\pi}4 & \frac{\pi}3 & \frac{\pi}2 & \frac{2\pi}3 & \frac{3\pi}4 & \frac{5\pi}6 & \pi
 \\[.3em] \hline&&&&&&&&& \\[-.8em]
r & 0 & 0.5 & 0.7 &  0.9 &  1 & 0.9 & 0.7 &  0.5 & 0
\\[.3em]  \hline
\end{array}
$$
\\[-1em]

\centerline{
\includegraphics[width=4.5in]
{10-3-Sin-Theta-Sample.png} 
}
\vspace{.5em}

\noindent 
As we said, the angles $\pi\leq\theta\leq 2\pi$ give negative radius and re-plot the same points.

%The computer does the same kind of plotting, but with many more points,
%so why bother with our piddly hand sketches and point plotting?
%Because of a very great danger for anyone who uses mathematics.
%If you let the computer do the thinking, not just the calculating, 
%you are ready to accept any bizarre wrong answer 
%without any way to check it. Then one typo error will escalate
%until your scientific paper has to be retracted,
%your company's expenses are ten times what you predicted,
%your bridge collapses, your rocket crashes.
%Don't let it happen!
%Before accepting the computer's answer, 
%you must check the expected answer qualitatively against a story or sketch, 
%and quantitatively by plotting sample points.

From the sketch, we may guess this curve is a circle, which we verify by converting 
to an $xy$-equation, and simplifying by completing the square:
$$
r =\sin(\theta) 
\quad\Longrightarrow\quad
\sqrt{x^2+y^2} \ =\ \,\frac{y}{\sqrt{x^2+y^2}}
\quad\Longrightarrow\quad
x^2+y^2-y = 0
$$
$$
\quad\Longrightarrow\quad
x^2+y^2-2(\tfrac12)y +(\tfrac12)^2 = (\tfrac12)^2
\quad\Longrightarrow\quad
x^2+(y{-}\tfrac12)^2 = (\tfrac12)^2\,.
$$
Indeed, this is a circle of radius $\frac12$ centered at $(x,y)=(0,\frac12)$.

\newpage

\noindent
{\bf More sketching.} We sketch the curve:
$$
r = 1+\sin(2\theta)\,.
$$
This is more complicated, so instead of computing a table of $\theta$ and $r$ values, 
we start by drawing the function $r = f(\theta)=1+\sin(2\theta)$ in our usual way as a rectangular graph,
labeling the horizontal and vertical axes by $r$ and $\theta$ because
that is how we intend to draw them later in the polar graph.
\\[1em]
\centerline{
\includegraphics[width=4.5in]
{10-3-Sin-2Theta-Rect.png} 
}
\vspace{-1em}

\noindent
Even without precise values, we can sketch the polar graph
by adjusting the radius according to the heights of the rectangular graph
(dotted lines).\footnote{
Graphically, we crush the entire horizontal $\theta$-axis in the rectangular graph
to the origin in the polar graph, spreading out the radial lines like a fan.
}
The blue lobe is traced by $\theta\in[-\frac{\pi}4, \frac{3\pi}4]$;
then the green lobe is for $\theta\in[\frac{3\pi}4, \frac{7\pi}4]$.
\\[1em]
\centerline{
\includegraphics[width=2.7in]
{10-3-Sin-2Theta.png} 
}
\vspace{0em}

\noindent We have $r=0$ at $\theta=\frac{3\pi}4$, so the tangent line
through the origin has slope $\tan \frac{3\pi}4=-1$, i.e.~$y = -x$.
Converting the original curve to rectangular coordinates gives:
$$\hspace{-.5em}
r \ = \ 1+\sin(2\theta) \ =\ 1 + 2\sin(\theta)\cos(\theta)
\quad\Longrightarrow\quad
\sqrt{x^2{+}y^2} \ =\ 1 +2\tfrac{y}{\sqrt{x^2{+}y^2}}\tfrac{x}{\sqrt{x^2{+}y^2}}
 \ = \ \tfrac{(x{+}y)^2}{x^2{+}y^2},
$$
or $(x^2+y^2)^3=(x+y)^4$. The polar form is much more informative about the shape!
%Implicit differentiation gives: 
%\\[1em]
%{\sc example:} Sketch the curve $r = \sin(2\theta)$. 
%Repeating the above procedure, we get four lobes, 
%traced in the order indicated 
%as we turn from $\theta=0$ to $\theta=2\pi$,
%with lobes 2 and 4 traced with $r<0$.
%\\[1em]
%\centerline{
%\includegraphics[width=2in]
%{10-3-Sin-2Theta.png} 
%}
%\vspace{0em}





\end{document}

%%%%%%%%%%%  Picture  %%%%%%%%%%
\\[1em]
\centerline{
%\includegraphics[width=2in]
{1-8-Discont-iv.png}
\qquad \qquad 
%\includegraphics[width=2in]
{1-8-Discont-v.png} 
}
\vspace{0em}

\noindent
%%%%%%%%%%%%%%%%%%%%%%%%%

%%%%%%%%%%%  Table  %%%%%%%%%%
In fact, we get the following derivatives:
$$\begin{array}{|c||c|c|c|c|c|c|}
\hline &&&&&& \\[-.9em] 
f(x) & \sin(x) & \cos(x) & \tan(x)  & \sec(x) & \csc(x) & \cot(x)
\\[.2em] \hline &&&&&& \\[-.9em]
f'(x) &\cos(x) & -\sin(x) & \sec^2(x) & \tan(x)\sec(x) & -\cot(x)\csc(x) & -\csc^2(x)
\\[.1em] \hline
\end{array}$$
\\[1em]
%%%%%%%%%%%%%%%%%%%%%%%%%%



















